\chapter{Background}
\pagenumbering{arabic}\hspace{3mm}

Depression and other mental illnesses have long been known to exhibit significant gender disparities in their reported prevalence \parencite{astbury1999gender}. In the case of depression — including both depressive symptoms and Major Depressive Disorder (MDD) — research in the United States has shown that “females experience 1.5- to 3-fold higher rates than males beginning in early adolescence” \parencite[p. 165]{APA2013DSM5}. However, studies have also found that this gender gap is considerably smaller in many developing countries \parencite{Sartorius1983}. This discrepancy raises the question of whether systematic biases—such as those introduced through data collection methods (e.g., census formats) or through Western gender norms and stereotypes—may be contributing to the observed differences in depression prevalence.

Despite this possibility, the gender disparity in depression has often been accepted as fact within the field of clinical psychology and embedded into the “gold standard” of psychiatric diagnosis, the \textit{Diagnostic and Statistical Manual of Mental Disorders, Fifth Edition} (DSM-5) \parencite{huntExploringGoldStandardEvidence2018}. Given these potential underlying biases in how mental health data is framed and interpreted, it is reasonable to expect that machine learning (ML) models trained to predict depression may also exhibit gender-based biases.

In this project, we propose to audit a high-performing ML model developed by Mahdi Ravaghi \parencite{ravaghiS04E11MentalHealth2024}, which won first place in Season 4, Episode 11 of Kaggle’s \href{https://www.kaggle.com/competitions/playground-series-s4e11}{Playground Series Competition}, titled \textit{Exploring Mental Health Data} \parencite{playground-series-s4e11}. The competition tasked participants with predicting mental health outcomes using a structured dataset, and submissions were evaluated solely based on accuracy. Ravaghi’s ensemble-based model achieved the highest accuracy among more than 2,500 submissions.

To evaluate both the performance and fairness of this model, we conducted a series of post hoc analyses including accuracy benchmarking, group-specific fairness metrics, and model interpretability techniques. We used Fairlearn’s \texttt{MetricFrame} to assess disparities in predictive outcomes across sensitive groups, applied post-processing mitigation through threshold optimization to address identified inequities, and used LIME (Local Interpretable Model-agnostic Explanations) \parencite{ribeiro2016why} to interpret how feature contributions vary at the individual level. While gender was our initial focus as the sensitive attribute, exploratory analysis led us to shift attention to employment status — specifically, whether individuals were working professionals or students — as this distinction revealed more pronounced disparities in model performance. Together, these methods aim to reveal not only how well the model performs, but also whom it serves—and whom it may potentially overlook or misclassify.
